\documentclass{beamer}
\usepackage{amsmath}
\usepackage{amssymb}

% Choose a theme
\usetheme{Madrid}
\usecolortheme{default}
\usefonttheme{serif}

% --- CUSTOM FOOTER ---
% Redefine the footline to remove author/institute and keep title/page number
\makeatletter
\setbeamertemplate{footline}
{
  \leavevmode%
  \hbox{%
  \begin{beamercolorbox}[wd=.5\paperwidth,ht=2.25ex,dp=1ex,center]{title in head/foot}%
    \usebeamerfont{title in head/foot}\insertshorttitle
  \end{beamercolorbox}%
  \begin{beamercolorbox}[wd=.5\paperwidth,ht=2.25ex,dp=1ex,right]{date in head/foot}%
    \usebeamerfont{date in head/foot}\insertframenumber{} / \inserttotalframenumber\hspace*{2ex}
  \end{beamercolorbox}}%
  \vskip0pt%
}
\makeatother
% --- END CUSTOM FOOTER ---

\title[Lecture 15: Plane wave propagation]{Lecture 15: Plane wave propagation}
\author{David Al-Attar}
\institute{Michaelmas term 2024}
\date{}

\begin{document}

% --- INTRODUCTION ---

\begin{frame}
    \titlepage
\end{frame}

\begin{frame}{Outline and Motivation}
    We will now study the simplest type of elastic waves: plane waves.
    \begin{itemize}
        \item We will analyze wave propagation in a \textbf{homogeneous} whole space (constant material parameters).
        \pause
        \item We will show that up to three different types of plane waves can propagate in any given direction.
        \pause
        \item For the special case of \textbf{isotropic} materials, we will derive the familiar P- and S-waves.
        \pause
        \item We will then use perturbation theory to see how these waves are affected by slight anisotropy.
        \pause
        \item Finally, we will qualitatively discuss how plane waves interact with boundaries (reflection and conversion).
    \end{itemize}
\end{frame}

% --- CHRISTOFFEL EQUATION ---

\begin{frame}{The Setup: Plane Waves in a Homogeneous Medium}
    We start with the linearised equations of motion for a \textbf{homogeneous} body, where $\rho$ and $A_{ijkl}$ are constant:
    $$ \rho\frac{\partial^{2}u_{i}}{\partial t^{2}} - A_{ijkl}\frac{\partial^2 u_{k}}{\partial x_{j}\partial x_{l}}=0 $$
    \pause
    We look for \textbf{plane wave} solutions of the form:
    $$ u_{i}(\mathbf{x},t)=f(t-\mathbf{p}\cdot\mathbf{x})a_{i} $$
    \pause
    \begin{columns}
    \column{0.5\textwidth}
        \begin{block}{Slowness Vector ($\mathbf{p}$)}
            \begin{itemize}
                \item Direction of $\mathbf{p}$ is the propagation direction.
                \item Magnitude sets the phase speed, $c = 1/||\mathbf{p}||$.
            \end{itemize}
        \end{block}
    \column{0.5\textwidth}
        \begin{block}{Polarisation Vector ($\mathbf{a}$)}
            \begin{itemize}
                \item A unit vector describing the direction of particle motion.
            \end{itemize}
        \end{block}
    \end{columns}
\end{frame}

\begin{frame}{Deriving the Christoffel Equation}
    We substitute the plane wave form $u_{i}=f(t-\mathbf{p}\cdot\mathbf{x})a_{i}$ into the equation of motion.
    \begin{itemize}
        \item The time derivative term is:
        $$ \rho\frac{\partial^{2}u_{i}}{\partial t^{2}} = \rho f'' a_i $$
        \pause
        \item The spatial derivative term is:
        $$ A_{ijkl}\frac{\partial^{2}u_{k}}{\partial x_{j}\partial x_{l}} = A_{ijkl} (f'' p_j p_l a_k) $$
    \end{itemize}
    \pause
    Equating these and cancelling the common factor $f''$ gives:
    $$ \rho a_i = A_{ijkl} p_j p_l a_k $$
    \pause
    Rearranging this gives the fundamental condition for a plane wave to exist.
\end{frame}

\begin{frame}{The Christoffel Equation}
    \begin{block}{The Christoffel Matrix}
        We define the Christoffel matrix $\mathbf{\Gamma}(\mathbf{p})$ as:
        $$ \Gamma_{ik}(\mathbf{p})=\frac{1}{\rho}A_{ijkl}p_{j}p_{l} $$
    \end{block}
    \pause
    \begin{alertblock}{The Christoffel Equation}
        The condition for plane wave solutions can now be written compactly as:
        $$ \Gamma_{ik}a_k = a_i \quad \text{or} \quad [\mathbf{\Gamma}(\mathbf{p})-\mathbf{I}]\mathbf{a}=0 $$
        This is a homogeneous system of linear equations for the polarisation vector $\mathbf{a}$.
    \end{alertblock}
    \pause
    \begin{itemize}
        \item For a non-trivial solution ($\mathbf{a} \neq 0$) to exist, the determinant of the matrix must be zero:
        $$ \det[\mathbf{\Gamma}(\mathbf{p})-\mathbf{I}]=0 $$
        This equation defines the set of all possible slowness vectors, a surface known as the \textbf{slowness surface}.
    \end{itemize}
\end{frame}

% --- EIGENVALUE PROBLEM ---

\begin{frame}{Reduction to a Symmetric Eigenvalue Problem}
    We can rewrite the Christoffel equation in a more familiar form.
    \begin{itemize}
        \item Let the slowness vector be $\mathbf{p} = \frac{1}{c}\hat{\mathbf{p}}$, where $c$ is the phase speed and $\hat{\mathbf{p}}$ is the direction.
        \pause
        \item The Christoffel matrix has the property $\mathbf{\Gamma}(\lambda \mathbf{p}) = \lambda^2 \mathbf{\Gamma}(\mathbf{p})$. Using this, the Christoffel equation becomes:
        $$ \mathbf{\Gamma}(\hat{\mathbf{p}})\mathbf{a} = c^2 \mathbf{a} $$
    \end{itemize}
    \pause
    \begin{alertblock}{An Eigenvalue Problem}
        For a fixed propagation direction $\hat{\mathbf{p}}$, this is an eigenvalue problem where the eigenvalues are the squared phase speeds ($c^2$) and the eigenvectors are the polarisation vectors ($\mathbf{a}$).
    \end{alertblock}
    \pause
    \begin{block}{Symmetry}
        Using the hyperelastic symmetry ($A_{ijkl}=A_{klij}$), one can show that the Christoffel matrix $\mathbf{\Gamma}(\hat{\mathbf{p}})$ is always symmetric.
    \end{block}
\end{frame}

\begin{frame}{Qualitative Solution}
    Because $\mathbf{\Gamma}(\hat{\mathbf{p}})$ is a $3 \times 3$ symmetric matrix, we know the properties of its solutions from linear algebra.
    \begin{itemize}
        \item It has \textbf{three real eigenvalues}: $c_1^2, c_2^2, c_3^2$.
        \pause
        \item It has \textbf{three corresponding eigenvectors} that are mutually \textbf{orthogonal}: $\mathbf{a}_1, \mathbf{a}_2, \mathbf{a}_3$.
    \end{itemize}
    \pause
    \begin{alertblock}{Physical Interpretation}
        In any given propagation direction $\hat{\mathbf{p}}$, there are (up to) \textbf{three distinct types of plane waves}.
        \begin{itemize}
            \item Each wave type has its own phase speed ($c_1, c_2, c_3$).
            \item The particle motions (polarisations $\mathbf{a}_1, \mathbf{a}_2, \mathbf{a}_3$) of these three waves are mutually orthogonal.
        \end{itemize}
    \end{alertblock}
\end{frame}

% --- ISOTROPIC CASE ---

\begin{frame}{Isotropic Case: P- and S-Waves}
    For an isotropic material, we can solve the eigenvalue problem explicitly. The Christoffel matrix becomes:
    $$ \Gamma_{ik}(\hat{\mathbf{p}})=\frac{\lambda+\mu}{\rho}\hat{p}_{i}\hat{p}_{k}+\frac{\mu}{\rho}\delta_{ik} $$
    \pause
    We can immediately find the eigenvectors and eigenvalues.
    \begin{block}{P-waves (Longitudinal)}
        If we choose the polarisation vector parallel to the propagation direction, $\mathbf{a} = \hat{\mathbf{p}}$:
        $$ \mathbf{\Gamma}(\hat{\mathbf{p}})\hat{\mathbf{p}} = \left(\frac{\lambda+2\mu}{\rho}\right) \hat{\mathbf{p}} $$
        \begin{itemize}
            \item This is a longitudinal wave, called a \textbf{P-wave}.
            \item The phase speed is $c_P = \alpha = \sqrt{(\lambda+2\mu)/\rho}$.
        \end{itemize}
    \end{block}
\end{frame}

\begin{frame}{Isotropic Case: S-Waves and Summary}
    \begin{block}{S-waves (Transverse)}
        If we choose $\mathbf{a}$ to be any vector orthogonal to $\hat{\mathbf{p}}$ (i.e., $\mathbf{a} \cdot \hat{\mathbf{p}}=0$):
        $$ \mathbf{\Gamma}(\hat{\mathbf{p}})\mathbf{a} = \left(\frac{\mu}{\rho}\right) \mathbf{a} $$
        \begin{itemize}
            \item This is a transverse wave, called an \textbf{S-wave}.
            \item The phase speed is $c_S = \beta = \sqrt{\mu/\rho}$.
            \item Since there are two orthogonal directions transverse to $\hat{\mathbf{p}}$, this eigenvalue is \textbf{two-fold degenerate}.
        \end{itemize}
    \end{block}
    \pause
    \begin{alertblock}{Summary for Isotropic Media}
        \begin{itemize}
            \item One type of P-wave with speed $\alpha$.
            \item Two types of S-waves with the same speed $\beta$.
            \item Speeds are independent of propagation direction.
            \item We always have $\alpha > \beta > 0$ in a solid.
        \end{itemize}
    \end{alertblock}
\end{frame}

% --- PERTURBATION THEORY ---

\begin{frame}{Slightly Anisotropic Materials}
    In some parts of the Earth, the material is slightly anisotropic. We can analyze this using perturbation theory.
    \begin{itemize}
        \item We write the elastic tensor as an isotropic part plus a small anisotropic perturbation:
        $$ A_{ijkl} = A^{\text{iso}}_{ijkl} + s~\delta A_{ijkl} $$
        \pause
        \item This leads to a perturbed Christoffel matrix:
        $$ \mathbf{\Gamma}(\hat{\mathbf{p}}) = \mathbf{\Gamma}^{\text{iso}}(\hat{\mathbf{p}}) + s~\delta\mathbf{\Gamma}(\hat{\mathbf{p}}) $$
        \pause
        \item We then find the first-order corrections to the isotropic eigenvalues ($c^2$) and eigenvectors ($\mathbf{a}$).
    \end{itemize}
\end{frame}

\begin{frame}{Perturbation of P- and S-waves}
    \begin{block}{Quasi P-wave}
        The non-degenerate P-wave is perturbed slightly.
        \begin{itemize}
            \item The wave is now a \textbf{quasi P-wave} (qP).
            \item The polarisation is almost, but not exactly, parallel to $\hat{\mathbf{p}}$.
            \item The phase speed now depends on the propagation direction.
        \end{itemize}
    \end{block}
    \pause
    \begin{block}{Quasi S-waves and Shear-Wave Splitting}
        The two-fold S-wave degeneracy is the most interesting case.
        \begin{itemize}
            \item Perturbation theory shows that the degeneracy is almost always lifted. The two S-waves now have slightly different speeds.
            \item This phenomenon is called \textbf{shear-wave splitting}.
            \item The resulting waves are \textbf{quasi S-waves} (qS1, qS2) with nearly transverse, orthogonal polarisations.
        \end{itemize}
    \end{block}
\end{frame}

\begin{frame}{Geophysical Significance of Shear-Wave Splitting}
    \begin{exampleblock}{Probing Mantle Flow}
        \begin{itemize}
            \item While most of the Earth appears isotropic to seismic waves (due to random crystal orientations), some regions in the mantle exhibit slight anisotropy.
            \pause
            \item This is caused by the large-scale alignment of mineral grains due to convective flow.
            \pause
            \item By observing shear-wave splitting, seismologists can make inferences about the direction and nature of flow within the Earth's mantle.
        \end{itemize}
    \end{exampleblock}
\end{frame}

% --- REFLECTION/TRANSMISSION ---

\begin{frame}{Waves at an Interface: Snell's Law}
    Consider two different homogeneous half-spaces meeting at a planar interface.
    \begin{itemize}
        \item An incident plane wave will strike the boundary. To satisfy the boundary conditions (continuity of displacement and traction), reflected and transmitted waves must be generated.
        \pause
        \item For the boundary conditions to hold true everywhere on the interface, the phase of all waves must match.
    \end{itemize}
    \pause
    \begin{alertblock}{Snell's Law}
        This phase-matching condition requires that the component of the slowness vector parallel to the interface must be conserved for all incident, reflected, and transmitted waves.
        $$ p_{\text{horizontal}} = \text{constant} $$
    \end{alertblock}
\end{frame}

\begin{frame}{Wave Conversion}
    \begin{block}{Reflection, Transmission, and Conversion}
        \begin{itemize}
            \item A single incident wave (e.g., a P-wave) will generally generate a full set of reflected and transmitted waves.
            \pause
            \item For an isotropic solid-solid interface, an incident P-wave will produce a reflected P-wave, a reflected S-wave, a transmitted P-wave, and a transmitted S-wave.
            \pause
            \item The generation of waves of a different type (e.g., an S-wave from an incident P-wave) is known as \textbf{wave conversion}.
        \end{itemize}
    \end{block}
    \pause
    \begin{exampleblock}{Complexity of Seismograms}
        Wave conversion is a fundamental property of elasticity and is a major reason why observed seismograms are so complex, containing many more arrivals than just the direct P- and S-waves.
    \end{exampleblock}
\end{frame}

% --- SUMMARY ---

\begin{frame}{Summary: What You Need to Know}
    \begin{itemize}
        \item How to derive the \textbf{Christoffel equation} from the equations of motion for a plane wave.
        \pause
        \item That this reduces to a \textbf{symmetric eigenvalue problem}, implying up to three orthogonal wave types in any direction.
        \pause
        \item How to solve the Christoffel equation for an isotropic solid to find \textbf{P-waves} (longitudinal) and \textbf{S-waves} (transverse, degenerate).
        \pause
        \item The effect of slight anisotropy: P-waves become \textbf{quasi-P}, and S-waves exhibit \textbf{shear-wave splitting}.
        \pause
        \item The basics of wave interaction at a boundary: \textbf{Snell's Law} (conservation of horizontal slowness) and \textbf{wave conversion}.
    \end{itemize}
\end{frame}

\end{document}
